\chapter{Background and Related Work}

\section{Introduction}


\section{Mapping production and reservoir data}


\begin{table}[]
\resizebox{\textwidth}{!}{%
\begin{tabular}{|l|lll|}
\hline
Data   Categories & \multicolumn{3}{l|}{Production  Engineering Component}                                                                                              \\ \hline
Static Data &
  \multicolumn{1}{l|}{\begin{tabular}[c]{@{}l@{}}Wellbore   Data Completion\\    \\ Class A\end{tabular}} &
  \multicolumn{2}{l|}{\begin{tabular}[c]{p{8cm}}Tubing Size,\\ Artificial Lift used with   its relevant characteristics\\ pump type and its   characteristic curve\\ Perforation Depth\end{tabular}} \\ \hline
Depth Data &
  \multicolumn{1}{l|}{\begin{tabular}[c]{@{}l@{}}Fluid Properties\\    \\ Class B\end{tabular}} &
  \multicolumn{2}{l|}{\begin{tabular}[c]{p{8cm}}Produced Fluid Data,\\ Interfacial Tension\\ Formation Volume Factor\\Density\\  Viscosity\end{tabular}} \\ \hline
\multirow{2}{*}{Time Data} &
  \multicolumn{1}{l|}{\multirow{2}{*}{\begin{tabular}[c]{@{}l@{}}Operating   Parameters\\    \\ Class C\end{tabular}}} &
  \multicolumn{2}{l|}{\begin{tabular}[c]{p{8cm}}It depends on well   completion (Gas Lift, Electrical Submersible Pumps, Sucker Rod Pumps, Natural   Flow, Gas Nat. Flow, Plunger Lift…). These parameters could me controlled parameters, manipulated paramters or disturbance paramters\\   Example \ac{ESP}:\\ Pump head\\ Motor temperature\\Frequency\\ Voltage\\ …\end{tabular}} \\ \cline{3-4}  \hline
\end{tabular}%
}
\caption{}
\label{production data}
\end{table}


\begin{table}[]
\resizebox{\textwidth}{!}{%
\begin{tabular}{@{}|l|ll|@{}}
\toprule
Data Categories &
  \multicolumn{2}{l|}{Reservoir Engineering   Component} \\ \midrule
\multirow{4}{*}{\begin{tabular}[c]{@{}l@{}}Reservoir Characteristics Data\\    \\ Class A\end{tabular}} &
  \multicolumn{1}{l|}{Geophysical data} &
  \begin{tabular}[c]{@{}l@{}}Seismic survey data\\    Well log data\end{tabular} \\ \cmidrule(l){2-3} 
 &
  \multicolumn{1}{l|}{Petrophysical data} &
  \begin{tabular}[c]{@{}l@{}}Permeability distributions\\    \\ Porosity distributions\\    \\ Formation Depth\\    \\ Reservoir Pressure\\    \\ Reservoir Temperature\\    \\ Fluid contact\end{tabular} \\ \cmidrule(l){2-3} 
 &
  \multicolumn{1}{l|}{Fluid Properties} &
  \begin{tabular}[c]{@{}l@{}}Fluid Composition\\    \\ PVT data\end{tabular} \\ \cmidrule(l){2-3} 
 &
  \multicolumn{1}{l|}{Rock/Fluid interaction data} &
  \begin{tabular}[c]{@{}l@{}}Relative Permeability data\\    \\ Capillary Pressure data\end{tabular} \\ \midrule
\begin{tabular}[c]{@{}l@{}}Project Design Parameters\\    \\ Class B\end{tabular} &
  \multicolumn{2}{l|}{\begin{tabular}[c]{@{}l@{}}Injection/Production well specification\\    \\ Well pattern design\\    \\ Well spacing\\    \\ Well architecture design\\    \\ EOR (Enhanced Oil Recovery) design parameters\end{tabular}} \\ \midrule
\begin{tabular}[c]{@{}l@{}}Field responses data\\    \\ Class C\end{tabular} &
  \multicolumn{2}{l|}{\begin{tabular}[c]{@{}l@{}}Fluid Production data\\    \\ Pressure data\\    \\ Project economics\end{tabular}} \\ \bottomrule
\end{tabular}%
}
\caption{Reservoir Data Categorization}
\label{reservoir_data}
\end{table}

\section{Production engineering applications}
lllllllllllllllllllllllll
fdddddddddddddddd
\section{Reservoir engineering applications}
sdddddddddddddddddddddd
%ToDo: include other sections here.
%\section{Further Content}

\section{Concluding remarks}
